% ============================================================================
% AERIS Paper - Section 7: Discussion
% Target Journal: Applied Intelligence (APIN) - Springer
% Version: 2026-01-18 (Data Authenticity Corrected)
% ============================================================================

\section{Discussion}
\label{sec:discussion}

This section provides critical analysis of AERIS's positioning, acknowledges
limitations, and offers practical guidelines for protocol selection.

\subsection{Honest Limitations}

We explicitly acknowledge the following limitations of AERIS:

\subsubsection{Energy Consumption vs PEGASIS}

AERIS consumes approximately \textbf{twice the energy} of PEGASIS (82.1mJ
vs. 41.9mJ at 100 nodes). This substantial overhead stems from:

\begin{enumerate}
    \item \textbf{Hierarchical transmission}: Multiple hops through skeleton
    backbone and gateway coordination, compared to PEGASIS's direct neighbor
    transmission.

    \item \textbf{Gateway load concentration}: Gateway nodes handle aggregated
    traffic, consuming disproportionate energy despite load balancing.

    \item \textbf{Control overhead}: Skeleton and gateway maintenance requires
    periodic reconfiguration messages.
\end{enumerate}

\textbf{Implication}: For applications where energy is the primary constraint,
PEGASIS remains the optimal choice. AERIS should not be selected for
energy-harvesting deployments with limited power budgets.

\subsubsection{Small-Scale No Advantage}

For networks with fewer than 100 nodes, \textbf{all protocols achieve 100\% PDR}.
AERIS's hierarchical overhead does not provide benefits at small scale:

\begin{itemize}
    \item At 50--100 nodes: LEACH, PEGASIS, and AERIS all achieve 100\% PDR
    \item AERIS advantage emerges only beyond 200 nodes where LEACH PDR
    degrades to $<$100\%
\end{itemize}

\textbf{Implication}: AERIS is not recommended for small-scale deployments
($<$100 nodes) where simpler protocols suffice.

\subsubsection{Execution Time}

AERIS requires significantly longer execution time than PEGASIS (15.17s vs
0.11s at 500 nodes). This overhead:

\begin{itemize}
    \item Increases setup time for each communication round
    \item Consumes additional computational energy not captured in the
    transmission energy model
    \item May impact applications requiring rapid reconfiguration
\end{itemize}

\subsubsection{CAS Module Contribution}

The ablation study reveals that the CAS (Context-Adaptive Switching)
module shows negligible independent effect (Hedges' $g = 0.00$) on
overall protocol performance. The primary performance contribution
comes from the Gateway coordination mechanism ($g = 10.09$).

\textbf{Implication}: CAS functions as a supporting framework rather
than the core innovation. Future work should focus on enhancing the
Gateway and skeleton backbone mechanisms rather than CAS refinement.

\subsection{AERIS vs. ML-Based Approaches}

Table~\ref{tab:ml_comparison_discussion} compares AERIS with recent
ML-based routing protocols.

\begin{table}[htbp]
\centering
\caption{AERIS vs. ML-Based Routing Protocols}
\label{tab:ml_comparison_discussion}
\begin{tabular}{@{}lccccc@{}}
\toprule
Aspect & AERIS & MeFi & MADRL & LSTM-routing \\
\midrule
Decision Time & $<$10ms & $\sim$600ms & $\sim$500ms & 50--80ms \\
Memory & 23KB & 2MB & 3.5MB & 700KB \\
Training & 0h & 48h & 96h & 16h \\
Adaptability & Rule-based & Learned & Learned & Learned \\
Interpretability & High & Low & Low & Low \\
Hardware & TelosB & Server & Server & Edge GPU \\
\bottomrule
\end{tabular}
\end{table}

\textbf{AERIS Advantages over ML}:
\begin{enumerate}
    \item \textbf{Zero training}: Immediate deployment without learning phase
    \item \textbf{Deterministic behavior}: Predictable routing decisions
    \item \textbf{Interpretability}: Rule-based logic enables debugging
    \item \textbf{Commodity hardware}: Deployable on 10KB RAM nodes
\end{enumerate}

\textbf{ML Advantages over AERIS}:
\begin{enumerate}
    \item \textbf{Adaptability}: ML protocols can learn environment-specific
    optimizations
    \item \textbf{Complex environments}: Neural networks may capture
    non-linear channel dynamics
    \item \textbf{Long-term optimization}: RL approaches optimize for
    network lifetime
\end{enumerate}

\textbf{Positioning}: AERIS targets resource-constrained deployments
where ML overhead is prohibitive. For server-class gateways with abundant
resources, ML approaches may provide superior long-term performance.

\subsection{Protocol Selection Guidelines}

Based on our comprehensive analysis, we provide practical guidelines
for protocol selection.

\subsubsection{Application-Based Recommendations}

Table~\ref{tab:recommendations} maps application requirements to
recommended protocols.

\begin{table}[htbp]
\centering
\caption{Protocol Selection Guidelines by Application (Based on Measured Data)}
\label{tab:recommendations}
\begin{tabular}{@{}p{3.5cm}ccp{3cm}@{}}
\toprule
Application & PDR Req. & Scale & Recommended \\
\midrule
Critical infrastructure & 100\% & $>$200 & \textbf{AERIS/PEGASIS} \\
Large-scale monitoring & 100\% & $>$200 & \textbf{AERIS/PEGASIS} \\
Dynamic topology & 100\% & Any & \textbf{AERIS} \\
Energy-critical sensing & $>$99\% & Any & \textbf{PEGASIS} \\
Small-scale deployment & Any & $<$100 & \textbf{LEACH} \\
\bottomrule
\end{tabular}
\end{table}

\subsubsection{Decision Flowchart}

We recommend the following decision process:

\begin{enumerate}
    \item \textbf{Is energy the primary constraint?} $\rightarrow$ Use PEGASIS
    (50\% energy savings vs AERIS)

    \item \textbf{Is network scale $<$100 nodes?} $\rightarrow$ Use LEACH
    (simplest, all protocols achieve 100\% PDR)

    \item \textbf{Is 100\% PDR required at scale ($>$200 nodes)?} $\rightarrow$
    Use AERIS or PEGASIS

    \item \textbf{Is AERIS's 18.5\% energy advantage over LEACH valuable?}
    $\rightarrow$ Use \textbf{AERIS}
\end{enumerate}

\subsection{Why AERIS Provides Value}

The protocol design space analysis reveals AERIS's positioning:

\begin{itemize}
    \item \textbf{LEACH}: Simple but PDR degrades at scale (98.68\% at 500 nodes)
    \item \textbf{PEGASIS}: Best energy efficiency, maintains 100\% PDR
    \item \textbf{HEED}: Balanced but neither optimal energy nor optimal PDR
    \item \textbf{AERIS}: 100\% PDR + 18.5\% energy reduction vs LEACH
\end{itemize}

\textbf{AERIS Value Proposition}:
\begin{enumerate}
    \item Better PDR stability than LEACH at scale
    \item Better energy efficiency than LEACH (18.5\% reduction)
    \item Robustness under node churn and regional failure
\end{enumerate}

\textbf{AERIS Limitation}:
\begin{enumerate}
    \item Higher energy than PEGASIS ($\sim$2$\times$)
    \item Longer execution time than all baselines
\end{enumerate}

\subsection{Threats to Validity}

\subsubsection{Internal Validity}

\begin{itemize}
    \item \textbf{Simulation vs. real-world}: Results are based on
    simulation with CC2420 energy model. Real-world deployment may reveal
    additional hardware-specific factors.

    \item \textbf{Channel model simplification}: Log-normal shadowing
    may not capture all propagation phenomena (multipath fading,
    temporal correlation).
\end{itemize}

\subsubsection{External Validity}

\begin{itemize}
    \item \textbf{Topology dependency}: Results are based on uniform
    random deployment. Structured deployments (grid, corridor) may
    exhibit different relative performance.

    \item \textbf{Traffic pattern}: Experiments assume periodic data
    generation. Event-driven or bursty traffic may affect protocol
    performance differently.
\end{itemize}

\subsubsection{Construct Validity}

\begin{itemize}
    \item \textbf{Latency not measured}: End-to-end transmission latency
    was not directly measured in experiments. Future work should implement
    hop counting to enable latency measurement.

    \item \textbf{Energy model scope}: Energy model covers transmission
    and reception but not computational overhead. AERIS's higher
    computational complexity is not fully captured.
\end{itemize}

\subsection{Future Work}

Based on our findings, we identify the following research directions:

\subsubsection{Latency Measurement Implementation}

A critical gap in current experiments is the lack of actual latency
measurement. Future work should:
\begin{itemize}
    \item Implement hop counting in the simulation framework
    \item Measure actual end-to-end transmission delay
    \item Validate theoretical complexity analysis with real data
\end{itemize}

\subsubsection{Gateway Mechanism Enhancement}

Given the Gateway component's dominant contribution ($g = 10.09$), future
work should focus on:
\begin{itemize}
    \item Adaptive gateway count based on network conditions
    \item Energy-aware gateway rotation strategies
    \item Multi-tier gateway hierarchies for ultra-large networks
\end{itemize}

\subsubsection{Lightweight ML Integration}

While maintaining AERIS's lightweight nature, future work could explore:
\begin{itemize}
    \item TinyML models ($<$5KB) for link quality prediction
    \item Quantized neural networks for gateway selection
    \item On-device learning with minimal memory overhead
\end{itemize}

\subsubsection{Hardware Validation}

Deploy AERIS on real TelosB/CC2650 hardware to validate:
\begin{itemize}
    \item Actual decision time and computational energy
    \item Memory footprint in constrained environments
    \item Real-world PDR under actual channel conditions
\end{itemize}
